% GENERATED FILE. Edit canonical lessons, then run npm run generate:books.
% canonical-source-hash: fnv1a64:e3e20f4860922582
% canonical-lessons: GU-C107-gadlu, GU-C107-cadar, GU-C107-dhabalo, GU-C107-pagthiyu, GU-C107-culo

\chapter{Mattress, Bedsheet, Blanket, Step, Stove}
\label{ch:gu-mattress-bedsheet-blanket-step-stove}
\hlchaptermodality{107}

\begin{chapteropening}
\textbf{By the end of this chapter:} \emph{I can say a mattress, a bedsheet, a blanket, a step and a stove.}

Complete the last lesson of chapter 107: i can say a mattress, a bedsheet, a blanket, a step and a stove.
\end{chapteropening}

\section[gādlũ]{\gu{ગાદલું} (gādlũ) — a mattress}

\label{lesson:GU-C107-gadlu}



\begin{quote}
Before the new one: say the Gujarati for soap, then the Gujarati for a towel.
\end{quote}

\subsection*{You'll want to know: \gu{ગાદલું}}
\textbf{\gu{ગાદલું}} (\emph{gādlũ}) — \textquotedblleft{}a mattress\textquotedblright{}.

\begin{grammarlens}[title={What you've built}]
One more word for this chapter.
\end{grammarlens}

\subsection*{Your turn}
\begin{itemize}
\raggedright
  \item \emph{Say it:} \emph{gādlũ}
  \item \emph{Say it:} \emph{gādlũ}, once more
  \item \emph{Say it:} say \emph{gādlũ}
  \item \emph{Recall:} say the Gujarati for soap, then the Gujarati for a towel, then say \emph{gādlũ} again
\end{itemize}

\begin{tcolorbox}[breakable,colback=teal!4,colframe=teal!35!black,title={Before you move on}]
What does \gu{ગાદલું} mean? (\textquotedblleft{}A mattress\textquotedblright{}.) Say it once more.
\end{tcolorbox}
\section[cādar]{\gu{ચાદર} (cādar) — a bedsheet}

\label{lesson:GU-C107-cadar}



\begin{quote}
Before the new one: say the Gujarati for a towel, then the Gujarati for a mattress.
\end{quote}

\subsection*{You'll want to know: \gu{ચાદર}}
\textbf{\gu{ચાદર}} (\emph{cādar}) — \textquotedblleft{}a bedsheet\textquotedblright{}.

\begin{grammarlens}[title={What you've built}]
One more word for this chapter.
\end{grammarlens}

\subsection*{Your turn}
\begin{itemize}
\raggedright
  \item \emph{Say it:} \emph{cādar}
  \item \emph{Say it:} \emph{cādar}, once more
  \item \emph{Say it:} say \emph{cādar}
  \item \emph{Recall:} say the Gujarati for a towel, then the Gujarati for a mattress, then say \emph{cādar} again
\end{itemize}

\begin{tcolorbox}[breakable,colback=teal!4,colframe=teal!35!black,title={Before you move on}]
What does \gu{ચાદર} mean? (\textquotedblleft{}A bedsheet\textquotedblright{}.) Say it once more.
\end{tcolorbox}
\section[dhābaḷo]{\gu{ધાબળો} (dhābaḷo) — a blanket}

\label{lesson:GU-C107-dhabalo}



\begin{quote}
Before the new one: say the Gujarati for a mattress, then the Gujarati for a bedsheet.
\end{quote}

\subsection*{You'll want to know: \gu{ધાબળો}}
\textbf{\gu{ધાબળો}} (\emph{dhābaḷo}) — \textquotedblleft{}a blanket\textquotedblright{}.

\begin{grammarlens}[title={What you've built}]
One more word for this chapter.
\end{grammarlens}

\subsection*{Your turn}
\begin{itemize}
\raggedright
  \item \emph{Say it:} \emph{dhābaḷo}
  \item \emph{Say it:} \emph{dhābaḷo}, once more
  \item \emph{Say it:} say \emph{dhābaḷo}
  \item \emph{Recall:} say the Gujarati for a mattress, then the Gujarati for a bedsheet, then say \emph{dhābaḷo} again
\end{itemize}

\begin{tcolorbox}[breakable,colback=teal!4,colframe=teal!35!black,title={Before you move on}]
What does \gu{ધાબળો} mean? (\textquotedblleft{}A blanket\textquotedblright{}.) Say it once more.
\end{tcolorbox}
\section[pagthiyũ]{\gu{પગથિયું} (pagthiyũ) — a step, a stair}

\label{lesson:GU-C107-pagthiyu}



\begin{quote}
Before the new one: say the Gujarati for a bedsheet, then the Gujarati for a blanket.
\end{quote}

\subsection*{You'll want to know: \gu{પગથિયું}}
\textbf{\gu{પગથિયું}} (\emph{pagthiyũ}) — \textquotedblleft{}a step, a stair\textquotedblright{}.

\begin{grammarlens}[title={What you've built}]
One more word for this chapter.
\end{grammarlens}

\subsection*{Your turn}
\begin{itemize}
\raggedright
  \item \emph{Say it:} \emph{pagthiyũ}
  \item \emph{Say it:} \emph{pagthiyũ}, once more
  \item \emph{Say it:} say \emph{pagthiyũ}
  \item \emph{Recall:} say the Gujarati for a bedsheet, then the Gujarati for a blanket, then say \emph{pagthiyũ} again
\end{itemize}

\begin{tcolorbox}[breakable,colback=teal!4,colframe=teal!35!black,title={Before you move on}]
What does \gu{પગથિયું} mean? (\textquotedblleft{}A step, a stair\textquotedblright{}.) Say it once more.
\end{tcolorbox}
\section[cūlo]{\gu{ચૂલો} (cūlo) — a stove}

\label{lesson:GU-C107-culo}



\begin{quote}
Before the new one: say the Gujarati for a blanket, then the Gujarati for a step.
\end{quote}

\subsection*{You'll want to know: \gu{ચૂલો}}
\textbf{\gu{ચૂલો}} (\emph{cūlo}) — \textquotedblleft{}a stove\textquotedblright{}.

\begin{grammarlens}[title={What you've built}]
One more word for this chapter.
\end{grammarlens}

\subsection*{Your turn}
\begin{itemize}
\raggedright
  \item \emph{Say it:} \emph{cūlo}
  \item \emph{Say it:} \emph{cūlo}, once more
  \item \emph{Say it:} say \emph{cūlo}
  \item \emph{Recall:} say the Gujarati for a blanket, then the Gujarati for a step, then say \emph{cūlo} again
\end{itemize}

\begin{tcolorbox}[breakable,colback=teal!4,colframe=teal!35!black,title={Before you move on}]
What does \gu{ચૂલો} mean? (\textquotedblleft{}A stove\textquotedblright{}.) Say it once more.
\end{tcolorbox}
